\section{Further properties of compact linear operators}

6 if the formula for eta looks familiar it is because it matrix multiplication. The T here is in fact encoding an infinity by infinity matrix. Notice that the condition of the squares of $alpha_jk$ is necessary for the multiplication with the matrix by x to be well-defined (Cauchy-Schwartz). This exercise shows that such matrices are only powerful enough to encode compact operators. General infinite matrices are not a very popular tool in FA for this reason
7 in fact not all compact operators are matrices either!
8 No: $l^infty$ is a huge space! Much bigger than $l²$. Why?
10 Notice this is actually also an infinite matrix, it is actually an infinity-by-infinity diagonal matrix! These are sometimes called diagonal operators. It turns out this exercise is actually an iff. That is, if T is compact, then $lambda_n -> 0$. If you want to prove this:
Step 1) All continuous functionals on l² are of the form < ., z>
Step 2) $e_n$ converge weakly to 0, since for every z, $<e_n, z> -> 0$. This follows from Parseval's equality and the fact the series in Parseval's equality converges
Step 3) Now apply Theorem 8.1-7

\begin{exercise}{6}
fill
\end{exercise}
\begin{proof}
fill
\end{proof} 

\begin{exercise}{7}
fill
\end{exercise}
\begin{proof}
fill
\end{proof} 

\begin{exercise}{8}
fill
\end{exercise}
\begin{proof}
fill
\end{proof} 

\begin{exercise}{10}
fill
\end{exercise}
\begin{proof}
fill
\end{proof} 
